\pg{353} 79. O ente composto depende, immediata ou mediatamente, do ente simples. O ente \emph{composto} depende, como dissemos, de uma causa extrinseca, Esta causa --- ou é um ente \emph{simples}, --- ou é um ente \emph{composto}, Se é um ente \emph{simples}, está provada a conclusão. Se é um ente \emph{composto}, renovamos a pergunta, até encontrarmos o ente \emph{simples}; porque, se não existisse o ente \emph{simples}, não poderia existir nenhum ente \emph{composto}. Logo o ente \emph{composto} depende, immediata ou mediatamente, do ente \emph{simples} \footnote{O ente \emph{composto} depende do \emph{simples}, enquanto se resolve, em ultima analyse, em elementos \emph{simples}; e, tambem e sobretudo, emquanto exige uma causa distincta, que una as partes componentes e que deve encontrarse finalmente no Ente subsistente na sua simplicidade. S. Thomaz prova esta proposição, como tambem a antecedebte. com as seguintes palavras: « Omnis compositio indiget aliquo componente. Si enim compositio est, ex pluribus est. Quae autem secundum se sunt plura, in unum non conveniunt, nisi ab aliquo componente uniautur, Si igitur compositus esset Deus, haberet componentem ; non enim ipse seipsum componere posset, quia nihil] est causa sui ipsius, esset enim prius seipso, quod est impossibile» (\emph{C}. Gent., 1, 18), E mais abaixo diz: « Ante multitudinem oportet invenire unitatem. In omni autem composito est multitudo. Tgitur oportet id, quod est ante omnia, seilicet \textsc{Deum}. omni compositione carere » (\emph{Ibid}.).}.

S0. Corollarios. --- Dos principios expostos resulta:

L. Na ordem \emph{ontologica}, ou da \emph{realidade}, o \emph{simples} precede o \emph{composto}, porque este depende d'aquelle; mas, na ordem \emph{logica}, ou do \emph{conhecimento}, o \emph{composto} precede o \emph{simples}, porque definimos o \emph{simples} pelo \emph{composto}, que se apresenta primeiro a nossa intelligencia, e dizemos que o \emph{simples} é o ente destituido de partes, é o \emph{não-composto} (\textasciicircum{}).

"\footnote{. Devemos. fazer outras considerações, relativas- ao comnhecimento do \emph{simples} e do \emph{composto}. --- O \emph{composto} pode ser conhecido quanto a uma ou mais paries e desconhecido quanto á outras: d'ahi os conhecimentos imperfeitos, relativos, constituidos por verdades e erros. --- \emph{O simples}, sendo uma coisa indivisivel, ou é conhecido totalmente, ou é totalmente desconhecido. Todavia, como o que é \emph{simples} em si mesmo, pode ser por nós considerado como uma coisa \emph{composta} e cognoscivel por varios juizos, alguns d'esses juizos podem ser verdadeiros e outros podem ser falsos. Assim Deus pode ser conhecido como Providencia e desconhecido como Espirito purissimo, € vice-versa. Se, porém, o \emph{simples} é tambem \emph{tal} para a nossa intelligencia, como são tantas verdades \emph{geometrieas}, \emph{ete}., então \emph{elle} é Ou inteiramente conhecido ou inteiramente \emph{iguorado}, \emph{sem} a \emph{possibilidade de} um termo medio,\par Vol, 1 24}

